\documentclass[11pt,a4paper]{article}
\usepackage{amsmath, amssymb, amsfonts}
\usepackage[utf8]{inputenc}
\usepackage[T1]{fontenc}
\usepackage{geometry}
\geometry{margin=2.5cm}
\usepackage{enumitem}
\usepackage{booktabs}
\usepackage{hyperref}
\usepackage{xcolor}
\usepackage{longtable}

\title{ICNNA Architectural Refactoring\\Phase~2: Architectural Decision (Session~3)\\
\large Steps~C and~D --- Deep Analysis and Synthesis for the Bounded ICNNA Update\\
Meeting Minute}
\author{Felipe Orihuela-Espina \and Claude (Computational Research Architect)}
\date{13 April 2026}

\begin{document}
\maketitle

% ================================================================
% Document Log
% ================================================================
\noindent\textbf{Document Log:}
\begin{itemize}[nosep]
    \item 13-Apr-2026: FOE/Claude --- Document created. Records Phase~2 Steps~C (deep analysis) and~D (synthesis) for the bounded ICNNA update path in MATLAB.
\end{itemize}

\vspace{1em}
\hrule
\vspace{1em}

% ================================================================
\section{Session Overview}

\begin{description}[style=nextline]
    \item[Date:] 13 April 2026
    \item[Participants:] Felipe Orihuela-Espina (FOE), Claude Opus 4.6 (Computational Research Architect role)
    \item[Predecessors:] Phase~2a minute (\texttt{Phase2a\_minute.tex}, 12-Apr-2026); Phase~2b minute (\texttt{Phase2b\_minute.tex}, 13-Apr-2026).
    \item[Objective:] Complete Phase~2 Steps~C and~D for ICNNA, now scoped as a \emph{bounded update in MATLAB} running in parallel with Cantor.
    \item[Outcome:] Steps~C and~D completed. Nine design principles (P1--P9) agreed for the bounded update. Key simplifications recorded: all current GUIs to be deprecated; MATLAB R2021a minimum; flat architecture to replace tree architecture; test infrastructure as early priority.
    \item[Model used:] Claude Opus 4.6
    \item[Status:] Phase~2 Steps~A--D complete. Steps~E (transitional strategy), F (open risks), and G (documentation and \texttt{claude/} folder organization) pending for next session.
\end{description}

% ================================================================
\section{Reframing: ICNNA as Bounded Update, Not Replacement}

Following the Phase~2b decision to create Cantor as a separate successor project, Phase~2's remaining steps were re-scoped to analyse and plan ICNNA's continued development as a \textbf{bounded update in MATLAB}. This is a modest version of the original Option~1 (``update ICNNA''), distinct from the full-scale refactoring that Cantor now represents.

\textbf{Two-tier structure agreed} (from earlier in the session):

\begin{description}
    \item[Survival tier:] What ICNNA \emph{must} deliver to remain usable for current users until Cantor reaches feature parity. Non-negotiable.
    \item[Enhancement tier:] What would make ICNNA more capable but could also be deferred to Cantor. Done opportunistically.
\end{description}

\textbf{Operating assumption:} Act \emph{as if} ICNNA were not yet doomed. Cantor's 9+~month timeline is likely longer given research-paced development. ICNNA must be genuinely viable during this period.

% ================================================================
\section{Step~C: Deep Analysis of the Bounded Update Path}

\subsection{Strengths (S1--S7)}

\begin{enumerate}[label=\textbf{S\arabic*.}]
    \item Mature domain knowledge embedded in 18~years of code --- not throwaway.
    \item Partial modernization in place: $\sim$20~classes in \texttt{+icnna} namespace with established pattern.
    \item SNIRF support substantially built (\texttt{icnna.data.snirf} package and \texttt{icnna.op.snirf.*} functions).
    \item Small, known user base ($\sim$5--10) --- no backward-compatibility burden at scale.
    \item MATLAB is native to the fNIRS community --- no adoption friction.
    \item No learning curve for developer --- every hour is productive work.
    \item Value-class decision settled after v1.3.1 lesson --- no further architectural ambiguity.
\end{enumerate}

\subsection{Limits (L1--L5)}

All limits are \emph{accepted as constraints}, not problems to solve. They are why Cantor exists.

\begin{enumerate}[label=\textbf{L\arabic*.}]
    \item MATLAB deployment ceiling ($\sim$2--3\,GB runtime). ICNNA stays within MATLAB; no standalone target.
    \item GUI framework ceiling. App Designer acceptable \emph{only} if full code control is retained (as in current manual implementation). To be considered during eventual GUI redo.
    \item Pass-by-value performance ceiling for large 3D tensors. Stays --- lesson learned from v1.3.1 handle-class experiment.
    \item AI integration ceiling. No AI work in ICNNA --- Cantor scope.
    \item Vendor lock-in to MathWorks. Accepted.
\end{enumerate}

\subsection{Hidden Constraints (HC1--HC6)}

\begin{enumerate}[label=\textbf{HC\arabic*.}]
    \item \textbf{GUI coupling.} Originally: each modernized class breaks its GUI(s), doubling effort. \textbf{Neutralized} by decision to deprecate all current GUIs in the next version. API modernization now proceeds independently of GUI work.
    \item \textbf{Dual architecture friction.} During transition, code must handle both \texttt{oo/} and \texttt{+icnna} types via conditional branching. No obvious shortcut --- accepted as unavoidable.
    \item \textbf{No test infrastructure.} Modernization without tests is high-risk. Test infrastructure is an \textbf{early roadmap priority} (v1.4.2).
    \item \textbf{Modernization order.} Originally strict bottom-up (leaf $\to$ container). \textbf{Relaxed} by the flat architecture decision: modernized classes hold IDs not nested objects, so container classes do not need dual-type handling of children. Bottom-up ordering still preferred but less rigid.
    \item \textbf{Research-priority interrupts.} Historically 2--3~times over ICNNA's life (including ICAF$\to$ICNNA transition). Manageable, but roadmap must have flexibility.
    \item \textbf{\texttt{@rawData\_Snirf} is a special case.} Logically belongs to SNIRF/BIDS work, not general rawData subtype modernization. May jump queue.
\end{enumerate}

\subsection{Failure Modes (FM1--FM5)}

\begin{enumerate}[label=\textbf{FM\arabic*.}]
    \item \textbf{Stuck-in-the-middle.} Modernization stalls, leaving a worse dual-architecture state than either pure-legacy or pure-modern. \textbf{Risk accepted as unavoidable.} Mitigated by sub-releases that each leave codebase in usable state.
    \item \textbf{GUI breakage alienates users.} \textbf{Eliminated} by decision to deprecate all GUIs up front. ICNNA goes GUI-less for several versions.
    \item \textbf{Cantor time competition.} Every hour on ICNNA delays Cantor. Mitigated by strict scope control and the enhancement-tier opt-out.
    \item \textbf{Standards drift (BIDS/SNIRF).} Deferred as concern --- cross this bridge when it comes.
    \item \textbf{MATLAB version incompatibility.} Resolved: \textbf{R2021a is the minimum supported version.} To be declared in startup script and documentation.
\end{enumerate}

\subsection{The Flat-Architecture Decision}

During Step~C discussion, an important architectural point was clarified: in the new \texttt{+icnna} architecture, constructs such as \texttt{experimentalUnit}, \texttt{session}, and \texttt{group} are stored \textbf{flat} within \texttt{experiment}, referenced by ID via \texttt{identifiableObject}, not nested as trees.

Implications:
\begin{itemize}
    \item Modernization is not a 1:1 class translation --- it is also a redesign of entity relationships.
    \item Each class spec in Phase~3 must consider its new flat-architecture interface, not just port the old one.
    \item Container classes (e.g. \texttt{experiment}) do not need to handle dual old/new types of children --- they hold IDs.
    \item Strict bottom-up modernization order is relaxed (see HC4).
\end{itemize}

% ================================================================
\section{Step~D: Synthesis and Recommendation}

\subsection{Core Synthesis}

The bounded ICNNA update has four defining characteristics:
\begin{enumerate}
    \item \textbf{A data-model modernization, not a full product update.} With GUIs deprecated and no AI/performance/deployment work, scope reduces to core class hierarchy, operations, plotting, tests, and documentation.
    \item \textbf{An architectural redesign within MATLAB, not a port.} Flat architecture replaces tree architecture. Modernized classes are redesigned interfaces, not faithful translations.
    \item \textbf{Bounded by Cantor's existence.} Genuinely architectural work (pipelines, AI, new analysis methods, GUI polish) belongs to Cantor. ICNNA's update sustains research utility during Cantor's build-out.
    \item \textbf{A one-person, research-paced project with interrupt tolerance.} Roadmap must be modular --- pausable, resumable, and reprioritizable without leaving codebase unusable.
\end{enumerate}

\subsection{Design Principles (P1--P9)}

\begin{enumerate}[label=\textbf{P\arabic*.}]
    \item \textbf{Data model first, everything else second.} Core data path (\texttt{signal}, \texttt{structuredData}, \texttt{dataSource}, \texttt{session}, \texttt{experiment}, and \emph{Definition} siblings) modernized before any enhancement-tier work.
    \item \textbf{Test-backed modernization from day one.} Every class modernized from v1.4.2 onwards ships with tests. A class without tests is not ``done.''
    \item \textbf{Sub-releases of 1--3 sessions, each leaving the codebase usable.} No half-refactored releases. All modernized classes compile and pass tests; legacy classes still in use still work; users can pick up any release and do real work.
    \item \textbf{GUI deprecation is a clean cut, not a slow fade.} Deprecate all GUIs in one version (v1.4.2 or v1.5.0). Document ICNNA as command-line/script-only during transition.
    \item \textbf{Flat architecture is a design task, not a translation task.} Each class's new interface is \emph{designed} (briefly but explicitly) before being \emph{implemented}. Operational definition of this principle to be developed progressively in Phase~3.
    \item \textbf{Bottom-up ordering, relaxed by flat architecture.} Leaf classes before containers, but with less strictness than pure trees would require.
    \item \textbf{Survival tier is the commitment; enhancement tier is opportunistic.} Survival tier (core data path modernization, tests, BIDS/SNIRF, docs, packaging) is non-negotiable. Enhancement tier (signal processing, statistics, pipelines) is done if time permits, droppable to Cantor without guilt.
    \item \textbf{Research-interrupt tolerance.} Roadmap must survive 2--3~weeks of no ICNNA work without losing coherence. Clear handover notes at each sub-release.
    \item \textbf{Strict scope containment.} No scope creep. ``Nice-to-have'' items discovered during modernization are logged for Phase~3 revision, not added to the current sub-release.
\end{enumerate}

\subsection{What the Update Is Not}

\begin{itemize}
    \item Not a performance project --- pass-by-value stays; no parallel processing work.
    \item Not a GUI project in the short term --- GUI redo is a future, separate effort \emph{within} this plan but late in it, built on the finished modern architecture.
    \item Not an AI project --- all AI-assisted analysis goes to Cantor.
    \item Not a new-methods project in general --- major new methods go to Cantor. However, some new methods in \texttt{+icnna.op} may be added on demand.
    \item Not a cross-platform project --- same MATLAB support matrix; no new platform testing.
\end{itemize}

\subsection{The Endpoint at v2.0.0}

\begin{itemize}
    \item All classes in \texttt{+icnna} package namespace; \texttt{oo/} folder removed.
    \item Flat architecture based on \texttt{identifiableObject}.
    \item A new GUI included (developed late in the modernization plan).
    \item BIDS compatibility delivered.
    \item Full test coverage for core data model.
    \item SNIRF compatibility maintained.
    \item Documentation complete for all public API.
\end{itemize}

\textbf{Important clarification:} v2.0.0 marks the end of this \emph{modernization plan}, but not the end of ICNNA's journey. The door remains open for future features, releases, and sub-releases as needs arise.

% ================================================================
\section{Decisions Recorded}

\begin{table}[h]
\centering
\begin{tabular}{lp{10cm}}
\toprule
\textbf{Item} & \textbf{Decision} \\
\midrule
Scope & Bounded ICNNA update in MATLAB, parallel to Cantor. Two-tier: survival + enhancement. \\
GUI strategy & Deprecate all current GUIs in the next version. New GUI built late in the plan on the finished modern architecture. \\
MATLAB minimum & R2021a. Declare in startup script and documentation when enforced. \\
Architecture & Flat (IDs via \texttt{identifiableObject}), not tree. Modernization is redesign, not port. \\
Pass-by-value & Final. No further handle-class experiments. \\
Tests & Early priority (v1.4.2). Mandatory for every modernized class onwards. \\
Principles & P1--P9 adopted as the design principles for the bounded update. \\
Endpoint & v2.0.0 marks end of modernization; not end of ICNNA. Future features possible. \\
\bottomrule
\end{tabular}
\caption{Decisions from Phase~2 Steps~C and~D.}
\end{table}

% ================================================================
\section{Pending Items for Next Session}

\begin{enumerate}[label=\textbf{Step~\Alph*:},start=5]
    \item \textbf{Transitional strategy} --- how ICNNA transitions from current dual state to v2.0.0 endpoint, how it coexists with Cantor during Cantor's build-out, and how the eventual handover to Cantor is managed.
    \item \textbf{Open risks / unknowns} --- explicit registry of what remains uncertain at the end of Phase~2.
    \item \textbf{Session documentation and \texttt{claude/} folder organization} --- including a future-proof scalable organization of the \texttt{claude/} folder itself.
\end{enumerate}

Phase~3 (roadmap definition) follows Phase~2 completion. FOE has prepared an initial prompt for Phase~3.

% ================================================================
\section{Documents Referenced}

\begin{enumerate}
    \item Phase~1 minute: \texttt{claude/Phase1\_minute.tex}
    \item Phase~2a minute: \texttt{claude/Phase2a\_minute.tex}
    \item Phase~2b minute: \texttt{claude/Phase2b\_minute.tex}
    \item Cantor project: \texttt{C:\textbackslash Users\textbackslash orihuelf-admin\textbackslash OneDrive\textbackslash Git\textbackslash Cantor}
    \item Global CLAUDE.md and project CLAUDE.md
    \item Memory files: user profile, ICNNA refactoring, Phase~2 Step~C outcomes, key documents
\end{enumerate}

\end{document}
